% Generated from ../chapters/22-from-experience-to-evidence.md
\chapter{From Experience to Evidence}

A strong personal observation deserves stronger measurement, not stronger rhetoric.

\section{Measurement hierarchy}

\subsection{Subjective}

Pain, ease, vitality, sleep quality, bodily awareness, and perceived symmetry matter because health is lived. They are also vulnerable to expectation and memory.

\subsection{Functional}

Strength, range of motion, balance, gait, reaction time, visual function, work capacity, and task performance are more discriminating.

\subsection{Structural}

Standardized photographs, 3D scans, dental measurements, ultrasound, DEXA, MRI, or CT may detect physical change. Imaging should be clinically justified; radiation and incidental findings have costs.

\subsection{Physiological}

Heart-rate variability, blood pressure, glucose regulation, cardiorespiratory fitness, inflammatory markers, and sleep measurements provide context, not a single score of youth.

\subsection{Molecular}

Epigenetic clocks, proteomics, metabolomics, immune profiling, and senescence-related markers are research tools with interpretation limits.

\section{A practical N-of-1 design}

\begin{enumerate}
\def\labelenumi{\arabic{enumi}.}
\tightlist
\item
  Write the hypothesis before reviewing the next result.
\item
  Select outcomes that could disconfirm it.
\item
  Standardize timing, equipment, posture, lighting, and effort.
\item
  Collect a baseline long enough to estimate ordinary variation.
\item
  Change one major factor where feasible.
\item
  Include withdrawal or lower-dose phases when safe.
\item
  Use blinded analysis for images or movement when possible.
\item
  Preserve raw data.
\item
  Report negative and ambiguous findings.
\item
  Distinguish statistical change from meaningful function.
\end{enumerate}

\section{Acoustic recordings}

Audio of neck, head, or mandibular events can document frequency and timing. It cannot identify anatomy on its own.

Useful analysis may include:

\begin{itemize}
\tightlist
\item
  event count and amplitude;
\item
  synchronization with motion capture;
\item
  association with range of motion or symptoms;
\item
  changes across months;
\item
  expert assessment where indicated.
\end{itemize}

\section{The ultimate test}

The framework predicts not merely improvement but \textbf{improved improvability}.

After one successful cycle, does the organism adapt faster, tolerate a broader range, or recover more completely from the next comparable challenge?

That second-order outcome is central to the escape-velocity hypothesis.
